\section{Introduction}
\label{sec:introduction}

Calligraphy---the writing of Chinese characters with a brush---is a
foundational form of traditional Chinese culture. Its principal script styles,
seal (\emph{zhuan}), clerical (\emph{li}), regular (\emph{kai}), running
(\emph{xing}), and cursive (\emph{cao}), employ distinct stroke techniques and
convey a wide range of qualities, from dignified and elegant to vigorous,
unrestrained, archaic, and lively. Two handwritten instances of the same
Chinese character can be equally
legible yet differ in the local structure relevant to calligraphy
(\emph{shufa}) instruction. A teacher may identify a shortened horizontal
stroke, a displaced hook, a low component, or an endpoint that fails to meet
its neighbour. Figure~\ref{fig:pipeline}(a) illustrates structural differences
using a regular-script reference and a running-script candidate of the same
character. Both are recognizable, but their stroke proportions, junctions,
and component spacing differ. These cross-style differences illustrate the
comparison task; they do not establish that either instance is incorrect.
Useful feedback must identify the relevant discrepancies and their locations
relative to the selected reference.

Global placement can vary without changing the local structure being
assessed. A character may be slightly larger, shifted on the canvas, or
tilted by a degree or two. Missing terminal fragments, over-wide strokes,
and misproportioned components should remain visible after such differences
are compensated. Reference-conditioned assessment must therefore separate
global placement from local discrepancy and return inspectable evidence
alongside a scalar score.

Recognition, style classification, and generation address character
identity, style category, and glyph synthesis, respectively
\cite{liu2013casia,zhang2026callinet,xu2026unicalli}. These objectives do not
directly evaluate spatially resolved, direction-specific discrepancies
between two instances. Formulating same-character structural comparison as
a document-analysis task exposes four computer-vision challenges.
\begin{enumerate}
  \item[C1] \emph{Overlapping labels under scarce annotation.} At crossings
  and adhesions, a pixel can carry multiple direction labels, which a
  mutually exclusive partition cannot represent. Endpoint evidence is sparse,
  and dense annotation is limited to 769 quality-controlled observations of
  40 character identities, without reliable writer identities.
  \item[C2] \emph{Global nuisance versus local structure.} Direct overlap
  confounds structure with translation, scale, and rotation. Flexible
  registration can erase proportion or component-placement errors, while an
  overlap score can reward direction channels that are empty in both images.
  \item[C3] \emph{Stability is not validity.} A score that is stable under
  allowed transforms may still disagree with how trained readers perceive
  structural similarity.
  \item[C4] \emph{A scalar does not localize.} Feedback must identify the
  direction family, the affected region, and whether ink is missing or extra.
  Fluent text cannot make an incorrect structural attribution reliable.
\end{enumerate}

\begin{figure*}[!t]
  \centering
  \includegraphics[width=\textwidth]{figures/redrawn/figure1_pipeline.pdf}
  \caption{Reference-conditioned structural assessment on one same-character
  pair. (a)~A regular-script reference and a running-script candidate of the
  same character are recognizable but structurally different. (b)~The candidate
  is parsed into overlapping direction channels vec1--vec5; the reference
  follows the same procedure. Black pixels carry multiple direction labels,
  and cyan outlines mark endpoints. The framed inset shows the multi-label
  bounding box at reduced scale. (c)~A bounded global transform registers the
  reference to the candidate, compensating global size and tilt (here scale
  1.20 and rotation \(3^\circ\)). The residual comprises overlap (dark grey),
  missing reference ink (blue, hatched), and extra candidate ink (red,
  counter-hatched). (d)~The aligned silhouette on the
  \(3\times3\) grid used for localized evidence; the values beneath it are
  the three ASDS spatial-distribution similarities of
  Sect.~\ref{sec:scores}. (e)~The structured output includes the production
  score, weakest direction family and its Dice, ink IoU, endpoint F1,
  transform, and missing/extra regions. The reported direction Dice belongs
  to that family, rather than the five-channel mean. All evidence is
  available before text generation. The cross-style pair is illustrative;
  the human study uses
  natural same-character, different-instance pairs from the project corpus. Source images: Calli-Tongji open subset \cite{yunmo_jixin_2025},
  CC~BY-NC~4.0.}
  \label{fig:pipeline}
\end{figure*}

Figure~\ref{fig:pipeline} outlines the framework, from overlapping structural
parsing and bounded registration to scalar agreement and localized evidence.
Parsing addresses C1, while registration and the score-semantics audit
address C2. The human study and diagnostic benchmark evaluate C3 and C4,
respectively, through the following research questions.
\begin{description}
  \item[RQ1 (representation and generalization).] Can convolutional and
  transformer parsers trained on the quality-controlled corpus of
  Sect.~\ref{sec:data} recover overlapping direction masks and sparse
  endpoint masks? How much does performance decline for character identities
  absent from training?
  \item[RQ2 (registration and score behaviour).] Relative to no alignment,
  does a constrained similarity registration reduce sensitivity to nuisance
  transforms while preserving penalties for controlled structural changes?
  Does a wider search improve this trade-off? How do channel-coverage
  semantics affect the frozen score, and does it distinguish same-character
  pairs on real references?
  \item[RQ3 (human validity).] How strongly do the frozen production score,
  the coverage-aware score, and aligned spatial-distribution similarity
  (ASDS) associate with blinded ratings of same-character structural
  similarity? Does the ASDS scalar depend on six-channel parsing?
  \item[RQ4 (localized diagnosis).] Which localized diagnostic outputs are
  reliable on interventions with known ground truth, and why do they fail?
\end{description}

The study makes four contributions aligned with these research questions.
\begin{enumerate}
  \item For RQ1, we define five independently activated direction-family
  channels and an endpoint channel, without a per-character prior at
  inference. The representation differs from Stroke-Seg's stroke instances
  (Sect.~\ref{sec:related}). Semantic quality control supports a matched
  three-seed comparison of U-Net, DeepLabV3+, and SegFormer-B2 on standard
  and character-disjoint splits. Ranking depends on
  target and protocol: DeepLabV3+ leads in-corpus direction parsing,
  SegFormer-B2 unseen-character direction parsing, and U-Net strict endpoint
  localization.
  \item For RQ2, we evaluate constrained similarity registration using
  controlled nuisance and structural perturbations of 200 real references.
  Paired comparisons cover no alignment and the tested wider, coarser search, alongside an
  inactive-channel score audit. Constrained registration
  substantially reduces nuisance penalties relative to no alignment, while
  structural edits remain penalized. Residual sensitivity and clipping-conditional
  summaries delimit this result.
  \item For RQ3, we compare the three scores in a blinded three-rater study
  and add a paired direct-ink control. ASDS development is retrospective,
  with character-grouped internal validation. Its observed association is
  strongest, but the direct-ink control supports a silhouette interpretation;
  direction-level evidence provides the rationale for the parser.
  \item For RQ4, we evaluate localized structured evidence and classify its
  failures. Required findings are usually surfaced, but exact-cell
  localization remains unreliable because of multi-cell truth and direction
  attribution errors. Language generation remains downstream of the
  structured evidence.
\end{enumerate}

The intended input is a raster exported from a direct digital writing canvas,
not a photograph. The claims concern reference-conditioned structural
analysis in this modality. They do not extend to camera robustness, writer-disjoint generalization,
stroke-order recovery, historical authenticity, or universal aesthetic
grading.
